\section{Risikomanagement-Integration je Phase}

Um eine terminologische Abgrenzung gegenüber dem aus der Softwareentwicklung 
bekannten Quality-Gate-Konzept zu schaffen, werden die phasenübergreifenden 
Prüfpunkte im Rahmen dieser Arbeit als ,,Risk Gates'' bezeichnet. 
Risk Gates werden nach jeder Hauptphase des DASC-PM platziert und bilden 
damit die formale Voraussetzung für den Übergang in die jeweils nachfolgende Phase.

Innerhalb jeder Hauptphase werden zunächst phasenspezifische Risiken 
identifiziert und bewertet, um die für die Entschiedungsfindung im Risk Gate
notwendigen Indikatoren zu erzeugen. Wie diese Indikatoren im Detail
aufgebaut sind, wird bewusst nicht vorgegeben, um den Einsatz dieser 
Strategie so individuell wie möglich zu gestalten.

Die Bewertung der Risiken geschieht anhand der in Kapitel \hyperref[sec:methode1]{\ref{sec:methode1}} vorgestellten Methoden
zur Risikobeurteilung in Data-Science-Projekten. Da die Konzentration relevanter
Risikotypen je nach Projektkontext und -ziel erheblich variiert, erfordert 
jede Phase eine individuell ausgerichtete Risikobewertung und ggf. eigens entwickelte 
Analyse und Bewertungsmethoden. Um den iterativen 
Charakter des DASC-PM zu erhalten und den Projektfortschritt nicht 
unverhältnismäßig zu verlangsamen, ist es dabei essenziell, dass die 
Risikobewertung von Beginn an in einen möglichst automatisierten Prozess 
überführt wird. Nur so können die Vorteile des Risk Gate Ansatzes mit jeder 
Iteration des Prozessmodells vollumfänglich genutzt werden. Dies kann und soll
die qualitative Bewertung von Risiken allerdings nicht ausschließen. Hierfür 
muss ein ebenfalls individueller Prozess erarbeitet werden.

Eine zentrale Bedingung für den Einsatz von Risk Gates ist, dass jede 
Risikobewertung ein dokumentiertes Ergebnisartefakt erzeugt, das mindestens 
einen messbaren KRI enthält. Diese Artefakte werden phasenübergreifend 
bereitgestellt und können in allen nachfolgenden Phasen referenziert werden. 
Dadurch entsteht ein kontinuierliches Wissensgedächtnis über den gesamten 
Projektverlauf: Erkenntnisse, die beispielsweise aus der Nutzungsphase 
hervorgehen, können genutzt werden, um den Projektantrag im Sinne einer 
iterativen Verbesserung nachzuschärfen, etwa um ethische Einschränkungen 
explizit aufzunehmen und gegebenenfalls unter Einbindung eines Fachbeirats 
zu konkretisieren.

\begin{figure}[H]
    \centering
    \includegraphics[width=0.85\textwidth]{RM-Extension_DASC-PM.jpeg}
    \caption{RM-Extension für das DASC-PM, eigene Darstellung nach \cite{schulz_dasc-pm_2022}}
    \label{fig:rm_extension}
\end{figure}

Die beschriebene Integration ist in Abbildung\hyperref[fig:rm_extension]{\ref{fig:rm_extension}} 
schematisch dargestellt. Das resultierende Modell wird als 
,,DASC-PM Risk Management Extension'', kurz ,,RM-Extension'', 
bezeichnet.

Zur Operationalisierung der Risk Gates wird ein dreistufiges Ampelsystem 
eingeführt, das den Risikostatus am Ende jeder Phase bewertet und die 
Entscheidung über den Phasenübergang strukturiert:

\begin{itemize}
    \item \textbf{Grün -- kein kritisches Risiko:} Alle Indikatoren liegen innerhalb 
    der definierten Schwellenwerte. Es sind keine zusätzlichen Maßnahmen 
    erforderlich. Der Übergang in die nächste Phase erfolgt ohne Einschränkungen.
    
    \item \textbf{Gelb -- erhöhtes Risiko:} Mindestens ein Indikatoren weist auf ein 
    erhöhtes, jedoch noch nicht kritisches Risiko hin. Das betreffende Risiko 
    wird als Artefakt dokumentiert, in allen nachfolgenden Gates weiterverfolgt 
    und kontinuierlich überwacht. Der Phasenübergang ist grundsätzlich möglich; 
    parallel können geeignete Gegenmaßnahmen eingeleitet werden.
    
    \item \textbf{Rot -- kritisches Risiko:} Mindestens ein Indikatoren überschreitet 
    den definierten kritischen Schwellenwert. Ein Übergang in die nächste Phase 
    ist erst zulässig, wenn geeignete Maßnahmen ergriffen wurden und der 
    betreffende Indikator mindestens in den Gelb-Bereich zurückgeführt wurde. 
    Das Risiko gilt damit als ausreichend behandelt, bleibt jedoch unter 
    kontinuierlicher Beobachtung. Die durchzuführende Maßnahme ist dabei
    phasenunspezifisch durchzuführen.
\end{itemize}

Es ist wichtig hervorzuheben, dass alle identifizierten Risiken 
kontinuierlich nachverfolgt werden müssen und somit in das ,,Risikogedächtnis'' einfließen. 
Der Unterschied zwischen den Ampelentscheidungen ist lediglich, dass 
bei einer grünen Ampel keine Maßnahme zu definieren, bei Gelb 
gibt es optionale Maßnahmen die parallel durchgeführt werden können 
und im roten Bereich ist eine Maßnahme unverzüglich durchzuführen.

Ein Risiko, das in einer frühen Phase, etwa der Datenerhebung, 
als erhöht (Gelb) eingestuft wurde, wird per  Carry-Forward-Logik 
in allen nachfolgenden Gates erneut anhand seiner Indikatoren geprüft. 
Verändert sich die Risikolage im weiteren Projektverlauf und überschreitet der Indikator den 
kritischen Schwellenwert, so wird der Status entsprechend der Überschreitung angehoben. 

Eine Ausnahme von der Regel, dass nach jeder Hauptphase ein Risk Gate folgt, 
bildet die Phase der ,,Nutzung''. Da das Modell zu diesem Zeitpunkt im 
produktiven Betrieb ist, greift anstelle eines einmaligen Gate-Prüfpunkts 
das in \hyperref[sec:risiko_monitoring]{Kapitel~\ref{sec:risiko_monitoring}} beschriebene kontinuierliche Monitoring. Dieses umfasst
sowohl technische Kennziffern als auch die fortlaufende Einschätzung des 
Risikoniveaus. Technische Indikatoren und Risikobeurteilung sind in dieser 
Phase konzeptionell eng miteinander verknüpft. Übersteigt ein Risiko im 
laufenden Betrieb einen kritischen Schwellenwert, so stehen in der Regel 
zwei strategische Optionen zur Wahl: die Weiterentwicklung des bestehenden 
Modells oder eine vollständige Neuentwicklung. Beide Szenarien münden in 
einem neuen Projektantrag und stoßen damit einen neuen Zyklus des DASC-PM an.